\section{Additional Econometric Details}
\label{app:econometrics}

This appendix records the implementation details that are fixed in the
production candidate universe. All quantities described below are recomputed
using only the training information admissible at the historical forecast
origin; the confirmatory evaluation sample is not used for hyperparameter
selection.

\subsection{Autoregressive and trailing-mean benchmarks}

The monthly inflation and labour AR(1) candidates estimate
\[
    y_\tau=\alpha+\phi y_{\tau-1}+\varepsilon_\tau
\]
by OLS on the admissible target history. The corresponding 12-month-mean
candidate is
\[
    \widehat y_{t}^{Mean12}
    =
    \frac{1}{12}\sum_{j=1}^{12} y_{t-j},
\]
provided the required target history is available. The GDP AR candidate uses
one lag with a constant and linear deterministic trend.

\subsection{Bridge Ridge implementation}

For each target and stage, the bridge dataset is rebuilt from the historical
information set before estimation. Predictors are standardised using moments
computed on the training sample, after which the Ridge objective in
Equation~\eqref{eq:ridge} is solved. The frozen penalties are
\[
    \lambda_{GDP}=10,\qquad
    \lambda_{\pi}=8,\qquad
    \lambda_{Labour}=10.
\]
The stage-specific feature matrix therefore changes as releases arrive, but
the penalty does not change with the confirmatory outcome.

For the GDP bridge system, transformed monthly indicators are aggregated to
quarterly predictors. When a current-quarter bridge feature has no current
quarterly observation, the most recently known quarterly feature value can be
carried forward under the frozen production builder; the affected feature is
recorded by the data audit. Inflation daily inputs are aggregated to monthly
frequency before transformation and may use the latest admissible value at or
before the required lagged month subject to a maximum carry of two months.
The analogous maximum carry in the labour builder is four months, with weekly
claims first aggregated to monthly means.

\subsection{Dynamic Factor Model}

The GDP DFM is implemented as a mixed-frequency state-space model using one
latent factor. The frozen configuration is:
\begin{itemize}[leftmargin=*]
    \item one factor and factor autoregressive order one;
    \item AR(1) idiosyncratic disturbances;
    \item standardisation of the observed panel;
    \item EM estimation with at most 100 iterations;
    \item convergence tolerance \(10^{-4}\).
\end{itemize}
The mixed-frequency input preserves ragged-edge missing observations and masks
the target-quarter GDP realization. Consequently, the DFM can condition on the
irregular release pattern without filling the panel with observations that
were not yet available.

\subsection{Factor Ridge implementation}

The labour Factor Ridge model standardises the admissible predictor panel and
extracts four static principal components. These four factors enter a Ridge
target regression with \(\lambda=8\). Both PCA and the penalised regression are
fit afresh using the admissible training sample at each forecast origin.

\subsection{Forecast combinations}

Two fixed equal-weight combinations enter the candidate universe. The
inflation Ridge--AR ensemble uses weights \((0.5,0.5)\) on Bridge Ridge and
AR(1). The labour equal-weight ensemble uses weights \((1/3,1/3,1/3)\) on
Bridge Ridge, AR(1), and Factor Ridge. The GDP Bridge--DFM ensemble uses fixed
weights \((0.5,0.5)\).

The rolling GDP Bridge--DFM ensemble instead estimates weights from prior
forecast performance. At origin \(t\), only forecast errors whose target
realizations were already observable may enter the weighting history. At most
the 12 most recent eligible errors are used, with a minimum history of 8.
Raw inverse-RMSE weights are bounded to the interval \([0.10,0.70]\) and
renormalised. If the minimum history is unavailable, the implementation uses
equal weights. These restrictions prevent current or future target outcomes
from affecting the combination assigned to a historical origin.

\subsection{Frozen policy-selection rule}

Stage-specific and fixed-comparator choices are made on common candidate
samples in the policy-development period. Within each target--stage cell,
models are ordered by development MSE, then MAE, then model name, with numerical
ties defined by tolerance \(10^{-12}\). For the target-level fixed comparator,
the same rule is applied after pooling the common development samples over all
stages. The exact resulting map is reported in Appendix~\ref{app:policies}.
